\documentclass[twocolumn, landscape]{article}
\usepackage[margin=1cm]{geometry}

\setlength{\columnseprule}{1pt}
\setlength{\parindent}{0pt}
\pagenumbering{gobble}

\usepackage{amsmath}
\usepackage{amsfonts}
\usepackage{amssymb}
\usepackage{amsthm}
\usepackage{thmtools}

\usepackage{unicode-math}
\setmainfont{TeX Gyre Pagella}
\setmathfont{TeX Gyre Pagella Math}

\usepackage{stmaryrd}
\usepackage{mathrsfs}
\usepackage{wasysym}

\usepackage[french]{babel}

\newtheoremstyle{myplain}%
  {\topsep}% Space above
  {\topsep}% Space below
  {\itshape}% Body font
  {}% Indent amount
  {}% Head font
  {}% Punctuation after head
  {5pt plus 1pt minus 1pt}% Space after head
  {\bfseries\thmname{#1}\thmnumber{ #2}.{\normalfont\thmnote{ (#3)}}} % Head spec

\newtheoremstyle{mydefinition}%
  {0.4\topsep}% Space above
  {0.4\topsep}% Space below
  {\normalfont}% Body font
  {}% Indent amount
  {}% Head font
  {}% Punctuation after head
  {5pt plus 1pt minus 1pt}% Space after head
  {\bfseries\thmname{#1}\thmnumber{ #2}.{\normalfont\thmnote{ (#3)}}} % Head spec


\declaretheorem[name=Théorème, style=mydefinition]{theo}
\declaretheorem[name=Définition, numberlike=theo, style=mydefinition]{defi}
\declaretheorem[name=Proposition, numberlike=theo, style=mydefinition]{prop}
\declaretheorem[name=Propriété, numberlike=theo, style=mydefinition]{prpr}
\declaretheorem[name=Conjecture, numberlike=theo, style=mydefinition]{conj}
\declaretheorem[name=Corollaire, numberlike=theo, style=mydefinition]{coro}
\declaretheorem[name=Lemme, numberlike=theo, style=mydefinition]{lemm}
\declaretheorem[name=Application, numberlike=theo, style=mydefinition]{appl}

\newtheoremstyle{myremark}%
  {0.5\topsep}% Space above
  {0.5\topsep}% Space below
  {\normalfont}% Body font
  {}% Indent amount
  {}% Head font
  {}% Punctuation after head
  {5pt plus 1pt minus 1pt}% Space after head
  {\itshape\thmname{#1}\thmnumber{ #2}.{\normalfont\thmnote{ (#3)}}} % Head spec
\declaretheorem[style=mydefinition,name=Remarque, numberlike=theo]{rema}
\declaretheorem[style=mydefinition,name=Exemple, numberlike=theo]{exem}
\declaretheorem[style=mydefinition,name=Convention, numberlike=theo]{conv}

\newcommand{\R}{\mathbb{R}}
\newcommand{\Rd}{\mathbb{R}^d}
\newcommand{\N}{\mathbb{N}}
\newcommand{\Z}{\mathbb{Z}}
\newcommand{\Q}{\mathbb{Q}}
\newcommand{\K}{\mathbb{K}}
\newcommand{\C}{\mathbb{C}}
\newcommand{\Ps}{\mathfrak{P}}
\newcommand{\dd}{\mathrm{d}}
\newcommand{\Tn}{\mathbb{T}^n}
\DeclareMathOperator{\imm}{im}
\DeclareMathOperator{\supp}{supp}
\DeclareMathOperator{\diam}{diam}
\DeclareMathOperator{\sgn}{sgn}
\DeclareMathOperator{\Aut}{Aut}
\DeclareMathOperator{\Gal}{Gal}
\DeclareMathOperator{\hess}{Hess}
\DeclareMathOperator{\Id}{Id}
\newcommand{\llb}{[\!\![}
\newcommand{\rrb}{]\!\!]}
\newcommand{\ol}[1]{\overline{#1}}
\newcommand{\ub}[2]{\underbrace{#1}_{#2}}
\newcommand{\wt}[1]{\widetilde{#1}}
\newcommand{\wh}[1]{\widehat{#1}}
\newcommand{\eqdef}{\stackrel{\mathrm{def}}{=}}
\newcommand{\wout}[1]{\setminus \{#1\}}
\newcommand{\abs}[1]{\left\lvert #1 \right\rvert}
\newcommand{\norm}[1]{\left\lVert #1 \right\rVert}
\newcommand{\normsub}[1]{\left\lVert \left\lvert #1 \right\rvert \right\rVert}

\AtBeginDocument{%à cause d'unicode-math
\let\leq\leqslant
\let\geq\geqslant
\let\epsilon\varepsilon
}

\usepackage[style=alphabetic]{biblatex}
\bibliography{../my}

\title{205. Espaces complets. Exemples et applications}

\begin{document}

\section{Espaces métriques complets}
Soit $(X, d)$ un espace métrique.

\subsection{Premières propriétés \cite{GouAna}}

\begin{defi}
    Une suite $(x_n) \in X^{\N}$ est dite \emph{de Cauchy} quand
    $$\forall \epsilon > 0, \exists N \geq 0, p \geq N \text{ et } q \geq N \implies d(x_p, x_q) < \epsilon.$$
\end{defi}

\begin{prop}
    Une suite convergente est de Cauchy. 
\end{prop}

\begin{exem}
    Soit $X= \Q$ muni de la distance usuelle. La suite $(x_n) \in \Q^{\N}$ définie par récurrence par $$x_0 = 1 \qquad x_{n+1} = \frac{x_n}{2} + \frac{1}{x_n}$$
    est de Cauchy mais ne converge pas.
\end{exem}

\begin{defi}
    L'espace métrique $X$ est dit \emph{complet} si toute suite de Cauchy $(x_n) \in X^{\N}$ converge dans $X$.
\end{defi}

\begin{prop}
    $\R$ est complet mais $\Q$ ne l'est pas. 
\end{prop}

\begin{prop}
    Soient $((X_i, d_i))_{i \in \{1, \dots, n\}}$ des espaces métriques. On définit sur $X= X_1 \times \dots \times X_n$ la distance produit 
    $$d((x_1, \dots, x_n), (y_1, \dots, y_n)) = \max_{i=1}^n d_i(x_i, y_i).$$
    Si $(X_i, d_i)$ est complet pour tout $i \in \{1, \dots, n\}$, alors $(X,d)$ est complet. 
\end{prop}

\begin{prop}
    Soit $d'$ une autre distance sur $X$. Si elle est équivalente à $d$, c'est-à-dire qu'il existe $c, C>0$ tel que pour tous $x,y \in X$, $$c d(x,y) \leq d'(x,y) \leq C d(x,y),$$
    alors $(X,d')$ est complet si et seulement si $(X,d)$ l'est. 
\end{prop}

\begin{coro}
    Tout espace vectoriel normé de dimension finie est complet pour la distance issue de sa norme.
\end{coro}

\begin{prop}
    Si $A \subset X$ est complet, alors $A$ est fermé dans $X$. À l'inverse si $X$ est complet et $A\subset X$ est fermé, alors $A$ est complet. 
\end{prop}

\begin{prop}
    Soient $(X, d_X)$ et $(Y, d_Y)$ deux espaces métriques avec $X$ non vide. On note $\mathscr{B}(X, Y)$ l'ensemble des fonctions $f:X\to Y$ bornées, c'est-à-dire telles que $f(X) \subset Y$ soit borné, et on définit
    $$d(f,g) = \sup_{x\in X} d_Y(f(x), g(x)).$$
    Cela fait de $\mathscr{B}(X,Y)$ un espace métrique, qui est complet si et seulement si $Y$ l'est. 
\end{prop}

\begin{prop}
    L'ensemble $\mathscr{C}_b(X,Y)$ des fonctions continues bornées est fermé dans $\mathscr{B}(X,Y)$. Il est donc complet si $Y$ l'est. 
\end{prop}

\begin{prop}
    Soit $Y$ un espace complet. Soit $A \subset X$ dense. Si $f : A \to Y$ est uniformément continue, il existe une unique fonction $g: X\to Y$ uniformément continue telle que $g|_A = f$.
\end{prop}

\begin{prop}
    Il existe un espace complet $\wh{X}$ et une isométrie injective $\iota : X \to \wh{X}$ avce $\iota{X}$ dense dans $\wh{H}$. Cet espace est appelé \emph{complété} de $X$, et il est unique à unique isométrie bijective près.  
\end{prop}

\begin{exem}
    $\R$ est le complété de $\Q$.
\end{exem}

\subsection{Théorème du point fixe et applications \cite{Demailly}}

\begin{defi}
    Une application $f: X\to X$ $k$-lipschitzienne pour $k<1$ est dite \emph{contractante}.
\end{defi}

\begin{theo}[Banach-Picard]
    Supposons $X$ complet. Alors toute application contractante $f : X\to X$ possède un unique point fixe.
\end{theo}

\begin{rema}
    La preuve du théorème montre que de plus la suite des itérées $$x_n = f^{(n)}(x_0)$$converge pour tout $x_0 \in X$ vers le point fixe, ce qui permet de le calculer de façon approchée.
\end{rema}

\begin{lemm}
    Soit $f: \mathscr{B}(x,r) \to \Rd$ tel que $f-\Id$ soit $k$-lipschitzienne pour $k<1$. Alors $f$ est un homéomorphisme de $\mathscr{B}(x,r)$ sur son image, et de plus $f^{-1}$ est $1/(1-k)$-lipschitzienne. On a enfin : 
    $$\mathscr{B}(f(x), (1-k)r) \subset f(\mathscr{B}(x,r)) \subset \mathscr{B}(f(x), (1+k)r).$$
\end{lemm}

\begin{appl}[inversion locale]
    Soit $U$ un ouvert de $\Rd$ et $f: U\to \Rd$ de classe $\mathscr{C}^1$. Supposons que $\dd_x f : \Rd \to \Rd$ soit bijective en $x\in \Rd$. Alors il existe un voisinage ouvert $V$ de $x$ dans $U$ et un voisinage ouvert $W$ de $f(x)$ tel que $f: V \to W$ soit un $\mathscr{C}^1$-difféomorphisme.
\end{appl}

\begin{lemm}
    Soit $f: C\to \Rd$ $k$-lipschitzienne avec $C = [t_0 -T, t_0 + T] \times \ol{\mathscr{B}}(y_0, r_0) \subset \R \times \Rd$. Soit $M = \sup_C \norm{f}$ et $T \leq \min(T_0, r_0/M)$. L'application
    $$\Phi : y \mapsto (t \mapsto \int_{t_0}^t f(s, y(s)) \,\dd s)$$
    laisse stable l'espace complet $\mathscr{C}([t_0 - T, t_0 + T], \ol{\mathscr{B}}(y_0, r_0))$ et est contractante.
\end{lemm}

\begin{appl}[Cauchy-Lipschitz]
    Soit $U \subset \R \times \Rd$. Soit $f: U \to \Rd$ localement lipschitzienne. Pour tout $(t_0, y_0) \in U$, le problème de Cauchy 
    $$\begin{cases}
        y'(t) = f(y(t)) \\
        y(t_0) = y_0
    \end{cases}$$
    possède une unique solution locale $y : [t_0 -T, t_0 + T] \to U$.
\end{appl}

\pagebreak

\subsection{Théorème de Baire et applications \cite{GouAna}}

\begin{theo}[Baire]
    Supposons $X$ complet. Soit $(\mathscr{O}_n)$ une suite d'ouverts denses de $X$. Alors $\bigcap O_n$ est encore dense, et en particulier non vide.
\end{theo}

\begin{coro}
    Supposons $X$ complet. Soit $(\mathscr{F}_n)$ une suite de fermés d'intérieur vide de $X$. Alors $\bigcup F_n$ est encore d'intérieur vide, et en particulier $\bigcup F_n \neq X$.
\end{coro}

\begin{appl}
    Il n'existe aucune norme faisant de $\R[X]$ un espace complet.
\end{appl}

\begin{appl}
    L'ensemble des fonctions $f: [0, 1] \to \R$ continues nulle part dérivable est dense dans $\mathscr{C}([0, 1], \R)$. En particulier, il existe une fonction réelle nulle part dérivable.
\end{appl}

\begin{appl}
    Soit $f : \R \to \R$ une fonction dérivable. Sa dérivée $f' : \R \to \R$ est continue sur un ensemble dense.
\end{appl}

\section{Espaces de Banach}

\subsection{Résulats fondamentaux \cite{Li}}

\begin{defi}
    Un espace de Banach est un espace vectoriel normé qui est complet pour la distance issue de sa norme.
\end{defi}

\begin{prop}
    Si $X$ est un espace métrique et $Y$ un espace de Banach, $\mathscr{B}(X,Y)$ et $\mathscr{C}_b(X,Y)$ sont des espaces de Banach. On définit aussi
    $$\mathscr{C}_0(X) = \{f\in \mathscr{C}(X, \R) \mid \forall \epsilon > 0, \exists K \subset X \text{ compact}, \sup_{x \not \in K} \abs{f(x)} < \epsilon\}.$$
    C'est un fermé de $\mathscr{C}_b(X, \R)$ donc un espace de Banach.  
\end{prop}

\begin{prop}
    Si $X$ et $Y$ sont deux espaces de Banach, alors $X\times Y$ et $\mathscr{L}_c(X, Y)$ aussi. En particulier, $X' = \mathscr{L}_c(X, \R)$, le dual (topologique) de $X$, est un espace de Banach.
\end{prop}

\begin{theo}[Riesz-Fischer]
    Soit $(\Omega, \mu)$ un espace mesuré. L'espace $L^p(\Omega, \mu)$ est complet pour tout $p \in [1, \infty]$.
\end{theo}

\begin{exem}
    En prenant pour $(\Omega, \mu)$ $\N$ muni de la mesure de comptage, on obtient que l'espace des suites
    $$\ell^p(\N) = \left\{ x_n \in \R^{\N} \mid \sum \abs{x_n}^p < \infty \right\}$$
    est complet. 
\end{exem}

\begin{theo}[Banach-Steinhaus]
    Soit $X$ un espace de Banach et $Y$ un espace vectoriel normé. Soit $(T_i)_{i\in I}$ une famille d'applications linéaires continues de $X$ dans $Y$. Alors :
    \begin{itemize}
        \item ou bien il existe $M < \infty$ tel que $$\forall i \in I, \normsub{T_i} \leq M$$
        \item ou bien l'ensemble des $x\in X$ tels que $$\sup \norm{T_i(x)}_Y = \infty$$
        est dense dans $X$. En particulier, il existe $x\in X$ tel que $$\sup \norm{T_i(x)}_Y = \infty.$$
    \end{itemize}
\end{theo}

\begin{theo}[application ouverte]
    Si $X$ et $Y$ sont deux espaces de Banach et que $T: X\to Y$ est linéaire continue bijective, alors $T^{-1} : Y \to X$ est aussi linéaire continue.
\end{theo}

\begin{theo}[graphe fermé]
    Soient $X$ et $Y$ deux espaces de Banach. Une application linéaire $T : X\to Y$ est continue si et seulement si son graphe $$\Gamma(T) = \{ (x, Tx) \in X\times Y \mid x\in X \}$$est fermé dans $X\times Y$.
\end{theo}

\subsection{Rappels d'analyse de Fourier \cite{Li}}

\subsection{Applications de la théorie des espaces de Banach à l'analyse de Fourier \cite{Rudin}}

\section{Espaces de Hilbert \cite{Li}}

\subsection{Projection sur un convexe fermé et conséquences}

\subsection{Compacité faible et applications}

\nocite{IPdev}
\printbibliography
\end{document}